
structure de @@M0@@-anneau sur un anneau @@M1@@, il revient au même de se donner un homomorphisme d'anneaux @@M2@@, tel que le diagramme 3.7.1 soit commutatif, ou un homomorphisme d'anneaux @@M3@@, tel que le diagramme 3.7.2 soit commutatif~; le premier définit les opérations @@M4@@, le second les opérations @@M5@@.
@@M6@@
@@M7@@

3.8. Soit @@M8@@ un anneau binômial (2.7)~; de (2.6.1) on déduit @@M9@@ les relations :
@@M10@@

Soit @@M11@@ un élément de @@M12@@~; la multiplication par @@M13@@ pour la loi @@M14@@ est caractérisée par sa valeur pour @@M15@@, soit :
@@M16@@
@@M17@@
En particulier, on voit que le @@M18@@-ième polynôme universel définissant la multiplication par @@M19@@ pour la loi @@M20@@ est le polynôme @@M21@@ (2.9).

